\subsection{Scheduling \texttt{Playbook} executions}
Kubernetes provides a built-in tool for scheduling jobs to run on a regular interval, which is called the \texttt{CronJob}.
A \texttt{CronJob} can be scheduled and will create a \texttt{Job} resource based on its \enquote{.spec.jobTemplate} field.
It also provides a history of a specified number of past executions and logs for these.
The \texttt{Job} itself then provides logs for the execution, dynamic resource scheduling and information on whether the \texttt{Job} was executed successfully.
Even though a \texttt{Job} is able to provide retry semantics, those should be disabled on created \texttt{Jobs}.
It is also important that the \enquote{.spec.concurrencyPolicy} field on the \texttt{CronJob} be set to \enquote{Forbid}, in order to avoid race conditions or collisions between \texttt{Playbook} executions.

\subsection{Runtime initialization}
Before the \texttt{Job} is able to execute the \texttt{Playbook} directly, it needs to be initialized. This includes cloning the repository with the \texttt{Playbook} and requirements file if applicable, as well as installing any dependencies specified in the requirements file.
An init-container is used in order to achieve these steps before executing the \texttt{Playbook}.
An init-container runs before application containers inside the \texttt{Pod} and may pass data to these containers, commonly via \enquote{EmptyDir} volumes.

\subsection{Executing \texttt{Playbooks}}
Once runtime initialization is complete, the application container can start and run the \enquote{ansible-playbook} command against the \texttt{Playbook} file, executing its contents.
The container needs to be set up in a way that it points at the correct files, including the \enquote{known\_hosts}, inventory and SSH keys.
